\documentclass[10pt,a4paper]{amsart}
\usepackage[margin=19mm]{geometry}
\usepackage{amsmath,amssymb,mathtools}
\usepackage[scr=rsfs,cal=euler]{mathalfa}
\usepackage{xfrac}
\usepackage[unicode,hidelinks,bookmarks=false]{hyperref}
\newcommand{\ie}{i.e.\ }
\newcommand{\cf}{cf.\ }
\newcommand{\resp}{respectively\ }
\newcommand{\Subsection}{\subsection}
\providecommand{\nobreakdash}{\nobreak\hbox{-}}
\setlength{\emergencystretch}{2em}
\multlinegap0pt
\usepackage{mathtools}
\usepackage{amssymb}
\usepackage{float}
\usepackage{enumerate}
\usepackage{tikz}
\usepackage[shortlabels]{enumitem}
\newtheorem{lem}{Lemma}
\def\dd{{\rm d}}
\title{NONLOCAL-TO-LOCAL LIMIT FOR LINEAR TRANSPORT EQUATIONS WITH MEASURE INITIAL DATA}
\author{Immanuel Ben Porat}
\date{Source: arXiv:2607.21968v1 (CC BY 4.0)}
\begin{document}
\maketitle

\noindent\emph{Source and licence.} NONLOCAL-TO-LOCAL LIMIT FOR LINEAR TRANSPORT EQUATIONS WITH MEASURE INITIAL DATA. Version arXiv:2607.21968v1 (CC BY 4.0), distributed by arXiv under the Creative Commons Attribution 4.0 International licence, http://creativecommons.org/licenses/by/4.0/, as recorded in the arXiv licence metadata for this exact version and shown on the abstract page https://arxiv.org/abs/2607.21968v1. Full licence notice: \texttt{licenses/arxiv-2607.21968-CC-BY-4.0.txt}.\par\smallskip

\noindent\textit{Editorial scope.} Counted unit: the article's lem principle (source0.tex lines 369--388). The article's complete original proof of it is delivered with it (source0.tex lines 389--447, one \texttt{proof} environment of 59 lines). No mathematics is written, repaired or re-derived: the statement and the proof are the article's own lines, in the article's own notation, and no retained line is substituted or re-flowed.  No further source-local context is needed: every cross-reference of every retained line resolves inside the two delivered spans.  Nothing was omitted from any retained span, so the package carries no omission record.  No retained line contains a citation, so the wrapper carries no bibliography and no bibliography entry.  No retained line contains a figure environment, an image-inclusion command or drawing code, so no image is delivered and the package declares no asset dependency.  Editorial formatting only, in addition: an AMS \texttt{amsart} preamble that restates the article's own declarations and macros used by the retained lines, namely the environments \texttt{lem} on separate counters, together with the standard packages those lines need.  The retained environments print in this document as Lemma 1 in order of appearance, whereas the article prints the same items with the numbering of its own full document.  The complete native source of the article as distributed in the arXiv e-print for this exact version is bundled unchanged as \texttt{sources/205902/source0.tex}.
\begin{lem}\label{Duhamel principle}
Suppose that: 
\begin{itemize}
    \item $W^{0}\in C^{\infty}(\mathbb{R}^{d};\mathbb{C})$.
    \item $A \in \mathrm{Lip}(\mathbb{R}^{d};\mathbb{R}^{d})\cap C^{\infty}(\mathbb{R}^{d};\mathbb{R}^{d})$.
    \item $S\in C^{\infty}(\mathbb{R}^{d};\mathbb{C})$ . 
\end{itemize}
\begin{enumerate}
    \item There exist a solution $\widetilde{W}\in W^{1,\infty}_{\mathrm{loc}}(\mathbb{R}_{+}\times \mathbb{R}^{d})$ to the  Cauchy problem 
\begin{align}
\partial_{t}\widetilde{W}-A\cdot \nabla_{x}\widetilde{W}=S+A\cdot \nabla_{x}W^{\mathrm{0}}, \ \widetilde{W}(0,\cdot)=0. \label{W tilde eq statement} 
\end{align}
\item If $\widetilde{W}$ is a solution of \eqref{W tilde eq statement} then   $W=W^{0}+\widetilde{W}$ is a solution of the Cauchy problem 
\begin{align}
\partial_{t}W-A\cdot \nabla_{x}W=S,\ W(0,\cdot)=W^{0}. \label{transport with source}    
\end{align}
\item  If $W_{1},W_{2}\in W^
{1,\infty}_{\mathrm{loc}}(\mathbb{R}_{+}\times \mathbb{R}^{d})$ are solutions to \eqref{transport with source} with the same initial data $W^{0}\in C^{\infty}(\mathbb{R}^{d};\mathbb{C})\cap L^{\infty}(\mathbb{R}^{d};\mathbb{C})$ and $S\in C^{\infty}(\mathbb{R}^{d};\mathbb{C})\cap L^{\infty}(\mathbb{R}^{d};\mathbb{C})$ then $W_{1}\equiv W_{2}$. 
\end{enumerate}
\end{lem}

\begin{proof}
\textbf{Step 1. Duhamel formula for $W^{0}=0$.} We consider the case where $W^{0}=0$. For any $s\in \mathbb{R}_{+}$ consider the IVP 
\begin{align}
\partial_{t}V_{s}-A\cdot \nabla_{x}V_{s}=0, \ V_{s}(s,\cdot)=S. \label{linear transport duhamel}    
\end{align}
By assumption $A\in \mathrm{Lip}(\mathbb{R}^{d};\mathbb{R}^{d})$, and therefore there exist a unique  solution $X(s,t,x)$ to the  ODE 
\begin{align}
\frac{\dd}{\dd t}X(s,t,x)=A(X(s,t,x)),\ X(s,s,x)=x.\label{Flow ODE}
\end{align} 
Moreover, we have that  $V_{s}(t,x)=X_{\ast}S$ is a solution to \eqref{linear transport duhamel}. We claim that 
\begin{align*}
W(t,x)\coloneqq \int_{0}^{t}V_{s}(t,x) \ \dd s     
\end{align*} 
is a solution to \eqref{W tilde eq statement}. By direct calculation it holds that   
\begin{align*}
\partial_{t}\left(\int_{0}^{t}V_{s}(t,x)\ \dd s\right)=V_{t}(t,x)+\int_{0}^{t}\partial_{t}V_{s}(t,x)\ \dd s= V_{t}(t,x)+\int_{0}^{t} A(x)\cdot \nabla_{x}V_{s}(t,x)\ \dd s 
\end{align*}   
and hence 
\begin{align*}
 \partial_{t}W(t,x)-A(x)\cdot \nabla_{x}W(t,x)=V_{t}(t,x)=S(x)   
\end{align*}
where in the last equation we used that  $X(t,t,x)=x$. \\
\textbf{Step 2}. We consider now the case of general $W^{0}\in C^{\infty}(\mathbb{R}^{d};\mathbb{C})$. By the previous step there exist a  solution $\widetilde{W}$  to the Cauchy problem 
\begin{align}
\partial_{t}\widetilde{W}-A\cdot \nabla_{x}\widetilde{W}=S+A\cdot \nabla_{x}W^{0}, \ \widetilde{W}(0,\cdot)=0. \label{W tilde eq}  
\end{align}
Define $W(t,x)\coloneqq W^{0}+\widetilde{W}$.  Then we have 
\begin{align*}
\partial_{t}W-A\cdot \nabla_{x}W&=\partial_{t}\widetilde{W}-A\cdot \nabla_{x}(\widetilde{W}+W^{0})\\
&=S+A\cdot \nabla_{x}(\widetilde{W}+W^{0})-A\cdot \nabla_{x}(\widetilde{W}+W^{0})=S.     
\end{align*}
It follows that $W=W^{0}+\widetilde{W}$ solves \eqref{transport with source} where $\widetilde{W}$ is a solution of \eqref{W tilde eq statement}.\\ 
\textbf{Step 3. Uniqueness.} If $W_{1},W_{2}$ are two solutions of \eqref{transport with source} with the same initial data $W^{0}$ then the difference $W_{1}-W_{2}$ satisfies  \begin{align}
\partial_{t}(W_
{1}-W_{2})-A\cdot \nabla_{x}(W_{1}-W_{2})=0, \ (W_{1}-W_{2})(0,\cdot)=0. \label{equation for difference W12}    
\end{align}
Since $A$ is real valued, by taking the real and imaginary parts of  \eqref{equation for difference W12} we conclude that 
\begin{align}
\partial_{t}\Re(W_{1}-W_{2})-A\cdot \nabla_{x}\Re(W_{1}-W_{2})=0 \label{real valued part} 
\end{align}
and 
\begin{align}
\partial_{t}\Im(W_{1}-W_{2})-A\cdot  \nabla_{x}\Im(W_{1}-W_{2})=0.\label{imaginary part}     
\end{align}
Multiplying \eqref{real valued part}-\eqref{imaginary part} by $\mathrm{sgn}(\Re(W_{1}-W_{2}))$ and $\mathrm{sgn}(\Im(W_{1}-W_{2}))$ respectively we get 
\begin{align}
\partial_{t}\left\vert \Re(W_{1}-W_{2})\right\vert -A\cdot \nabla_{x}\left\vert \Re(W_{1}-W_{2})  \right\vert=0 
\label{real absolute}
\end{align}
and 
\begin{align}
\partial_{t}\left\vert \Im(W_{1}-W_{2})\right\vert-A\cdot \nabla_{x}\left\vert \Im(W_{1}-W_{2})\right\vert=0. \label{imaginary absolute}      
\end{align}
Evaluating \eqref{real absolute}-\eqref{imaginary absolute} at a maximum point of $\left\vert \Re(W_{1}-W_{2})\right\vert$ and $\left\vert \Im(W_{1}-W_{2})\right\vert$ respectively we conclude that 
\begin{align*}
\frac{\dd}{\dd t}\left\Vert \Re(W_{1}-W_{2})(t,\cdot)\right\Vert_{\infty}=\frac{\dd}{\dd t}\left\Vert \Im(W_{1}-W_{2})(t,\cdot)\right\Vert_{\infty}=0,     
\end{align*}
so that $\Re(W_{1})=\Re(W_{2}), \ \Im(W_{1})=\Im(W_{2})$. Thus $W_{1}\equiv W_{2}$, proving uniqueness. 
\end{proof}

\end{document}
